\section{Geometric information about large systolic ratios}

Within the sampled product distributions, low normalized ridge area and
balanced symplectic vertex covariance each identify populations with
larger mean numerical systolic ratios. These give two different geometric
descriptions of promising bodies: one through their two-dimensional faces,
the other through their extreme points. Both respect affine symplectic
changes of coordinates and can be evaluated before capacity. The evidence
supports geometric selection; the square--HKO comparison above shows why
it cannot turn small ridge area into a universal rule for increasing the
ratio of an individual body.

The covariance condition makes the second description precise. For
vertices $v_1,\ldots,v_N$, put
\[
 \overline v=\frac1N\sum_{i=1}^Nv_i,
 \qquad \Gamma=\frac1N\sum_{i=1}^N(v_i-\overline v)(v_i-\overline v)^T.
\]
In suitable symplectic coordinates, this matrix has Williamson normal
form $\operatorname{diag}(\nu_1,\nu_2,\nu_1,\nu_2)$, with
$0<\nu_1\leq\nu_2$. The ratio $\rho(K)=\nu_2/\nu_1$
measures the imbalance between the covariance scales of the two
symplectic coordinate planes. It equals one when those scales agree.

This score is invariant under affine symplectic maps and common dilation.
Indeed, centering removes translations, and a symplectic map $S$ gives
$\Gamma(SK)=S\Gamma(K)S^T$. With
$J=\left(\begin{smallmatrix}0&I_2\\-I_2&0\end{smallmatrix}\right)$,
\[
 J\Gamma(SK)=S^{-T}\bigl(J\Gamma(K)\bigr)S^T.
\]
Thus the eigenvalues $\pm i\nu_1,\pm i\nu_2$ of $J\Gamma$ are unchanged.
A dilation multiplies both $\nu_j$ by the same squared scale factor,
preserving their ratio.

For a product, if $A$ and $B$ are the planar vertex covariance matrices
of its factors, then $\Gamma=\operatorname{diag}(A,B)$. Consequently
$(J\Gamma)^2=\operatorname{diag}(-BA,-AB)$, so $\nu_1^2,\nu_2^2$ are
the eigenvalues of $AB$.

\begin{samepage}
The following contrasts give the empirical support for these conditions.
Each compares selections within the same generated population and
factor-side-count groups; the bodies were selected from their geometry
before their ratios were evaluated.
\begin{center}
\begin{tabular}{lrr}
\toprule
Selection & Compared with & Mean difference\\
\midrule
Low $R$ & Comparison bodies & $+0.311$\\
Lower $C$ at low $R$ & Higher $C$ at low $R$ & $+0.034$\\
Low $\rho$ & Comparison bodies & $+0.333$\\
Low $\rho$ & Low $R$, then lower $C$ & $+0.015$\\
\bottomrule
\end{tabular}
\end{center}
\end{samepage}
The first contrast has selected and comparison means $0.626$ and $0.316$.
The second splits the lowest one percent by $R$ into halves according to
$C$; its difference is positive in eight of ten side-count groups. Thus
the dominance of the largest ridge adds information even when total ridge
area is already small.

The two covariance contrasts come from a separate comparison on two
fresh pools, motivated by an exploratory inverse association between
$\rho$ and the ratio. Their descriptive $95\%$ Student-$t$ intervals,
using the twenty pool-by-group differences, are $[0.310,0.356]$ and
$[-0.015,0.044]$. These contrasts support both selection criteria,
without resolving which yields larger mean ratios. The geometric reason
for their association with capacity remains open.

\subsection{Changing the population and strengthening the filter}

Ridge area also helped select generic ten-facet polytopes. Here the number
of ridges varies even though the number of facets is fixed. To compare the
area per ridge, we used $R$ divided by the ridge count. The lowest one
percent by this quantity had mean systolic ratio $0.626$, compared with
$0.339$ for comparison bodies from the same pool. More severe filtering
gave no further improvement in this sample: the lowest tenth of the
selected group had a mean ratio $0.032$ below the remaining selected bodies.
The descriptive $95\%$ interval from resampling these selected bodies
was $[-0.095,0.033]$, which spans zero. This interval conditions
on the realized selection; it does not include variation from generating
and selecting a new pool.

The product filters also remained useful when the factors were tangential
polygons. These polygons have every side tangent to a circle. We obtained
them by keeping the sampled unit normals and setting every support height
to one, retaining pairs for which both the original and the changed
polygons were valid. For products with side counts $4\times6$ and
$6\times6$, selection by small $\rho$ increased the mean ratio by $0.259$
relative to comparison bodies; ridge selection increased it by $0.225$.
These are averages of the two group differences in one pool. Together
with the original product experiment, they show that the usefulness of
both filters survived this particular change of sampling distribution.

\section{Changing a body or changing the sampling distribution}

Selection uses geometry to choose among sampled bodies. We can also ask
whether changing that geometry raises the ratio of an individual body,
or whether the successful samples suggest a better sampling distribution.

\subsection{Making the factors tangential}

Replacing both factors by tangential polygons sometimes raised and
sometimes lowered the systolic ratio. With the normals fixed, setting
all support heights to one gave a mean numerical increase of $0.015$
among the tested $4\times4$ and $4\times6$ products, with changes of both
signs. Thus a useful selection rule for tangential products does not
supply a rule for improving every product by making its factors tangential.

This comparison changes the support heights, then scales each factor to
area one. The latter normalization preserves the systolic ratio: for
$a,b>0$,
\[
 \operatorname{diag}(aI_2,bI_2)
 =\sqrt{ab}\,\operatorname{diag}(\sqrt{a/b}\,I_2,\sqrt{b/a}\,I_2),
\]
a common dilation followed by a symplectic map. The comparison concerns
products for which the original factors and both height replacements
remain admissible; it measures the change of shape within that family.

The tested ambient rotations $K\mapsto UK$, $U\in SO(4)$, also gave
smaller search maxima than fresh sampling at equal evaluation budgets
(Appendix~\ref{app:trial-ambient}).

\begin{samepage}
\subsection{Learning a sampling distribution}

Adapting the sampling distribution improved the search for pentagon
products in each of three comparisons. The cross-entropy
method~\cite{dsDeBoer2005} repeatedly fits a distribution to the best part
of the current population and draws a new population from it. Under the
same number of capacity evaluations, it attained both a larger maximum
systolic ratio and a larger median among the eight best distinct bodies:
\begin{center}
\begin{tabular}{crrrr}
\toprule
 & \multicolumn{2}{c}{Maximum} & \multicolumn{2}{c}{Median of best eight}\\
Run & Independent & Adaptive & Independent & Adaptive\\
\midrule
1 & $0.760$ & $0.799$ & $0.729$ & $0.782$\\
2 & $0.811$ & $0.831$ & $0.754$ & $0.776$\\
3 & $0.760$ & $0.856$ & $0.722$ & $0.814$\\
\bottomrule
\end{tabular}
\end{center}
\end{samepage}
The gain therefore extended beyond a single best observation in each run.

The sampler described each factor by log ratios of its normal-angle gaps
and centered logarithmic lengths of its dual normals, together with the
relative angle between the factors. These coordinates remove separate
factor scales and simultaneous planar rotation, which preserve the
systolic ratio. Independent Gaussian coordinates described gaps and
lengths; the relative angle used a wrapped Gaussian with a circular mean
update. Each update shifted the distribution's moments towards those of
the quarter of the population with largest ratios.

The comparison also includes the numerical evaluator's input restrictions.
It accepts only bodies within a prescribed coordinate range. Independent
sampling uses the original coordinates, whereas the adaptive method
reconstructs a normalized representation, so this restriction affects
the two distributions differently. The observed gain belongs to the
implemented combination of adaptation and admission. All recorded ratios
remained below one.
